\documentclass[10pt,a4paper]{amsart}
\usepackage{amsmath,amssymb,amsthm,graphicx}
\usepackage{verbatim}
\newenvironment{reference}{\comment}{\endcomment}
%\newenvironment{reference}{}{}
\newenvironment{slogan}{\comment}{\endcomment}
\newenvironment{history}{\comment}{\endcomment}

% For commutative diagrams we use Xy-pic
\usepackage[all]{xy}

% We use 2cell for 2-commutative diagrams.
\xyoption{2cell}
% Theorem environments.
%
\theoremstyle{plain}
\newtheorem{theorem}[subsection]{Theorem}
\newtheorem{proposition}[subsection]{Proposition}
\newtheorem{lemma}[subsection]{Lemma}

\theoremstyle{definition}
\newtheorem{definition}[subsection]{Definition}
\newtheorem{example}[subsection]{Example}
\newtheorem{exercise}[subsection]{Exercise}
\newtheorem{situation}[subsection]{Situation}

\theoremstyle{remark}
\newtheorem{remark}[subsection]{Remark}
\newtheorem{remarks}[subsection]{Remarks}

\numberwithin{equation}{subsection}

% Macros
%
\def\lim{\mathop{\mathrm{lim}}\nolimits}
\def\colim{\mathop{\mathrm{colim}}\nolimits}
\def\Spec{\mathop{\mathrm{Spec}}}
\def\Hom{\mathop{\mathrm{Hom}}\nolimits}
\def\Ext{\mathop{\mathrm{Ext}}\nolimits}
\def\SheafHom{\mathop{\mathcal{H}\!\mathit{om}}\nolimits}
\def\SheafExt{\mathop{\mathcal{E}\!\mathit{xt}}\nolimits}
\def\Sch{\mathit{Sch}}
\def\Mor{\mathop{\mathrm{Mor}}\nolimits}
\def\Ob{\mathop{\mathrm{Ob}}\nolimits}
\def\Sh{\mathop{\mathit{Sh}}\nolimits}
\def\NL{\mathop{N\!L}\nolimits}
\def\CH{\mathop{\mathrm{CH}}\nolimits}
\def\proetale{{pro\text{-}\acute{e}tale}}
\def\etale{{\acute{e}tale}}
\def\QCoh{\mathit{QCoh}}
\def\Ker{\mathop{\mathrm{Ker}}}
\def\Im{\mathop{\mathrm{Im}}}
\def\Coker{\mathop{\mathrm{Coker}}}
\def\Coim{\mathop{\mathrm{Coim}}}

% Boxtimes
%
\DeclareMathSymbol{\boxtimes}{\mathbin}{AMSa}{"02}

%
% Macros for moduli stacks/spaces
%
\def\QCohstack{\mathcal{QC}\!\mathit{oh}}
\def\Cohstack{\mathcal{C}\!\mathit{oh}}
\def\Spacesstack{\mathcal{S}\!\mathit{paces}}
\def\Quotfunctor{\mathrm{Quot}}
\def\Hilbfunctor{\mathrm{Hilb}}
\def\Curvesstack{\mathcal{C}\!\mathit{urves}}
\def\Polarizedstack{\mathcal{P}\!\mathit{olarized}}
\def\Complexesstack{\mathcal{C}\!\mathit{omplexes}}
% \Pic is the operator that assigns to X its picard group, usage \Pic(X)
% \Picardstack_{X/B} denotes the Picard stack of X over B
% \Picardfunctor_{X/B} denotes the Picard functor of X over B
\def\Pic{\mathop{\mathrm{Pic}}\nolimits}
\def\Picardstack{\mathcal{P}\!\mathit{ic}}
\def\Picardfunctor{\mathrm{Pic}}
\def\Deformationcategory{\mathcal{D}\!\mathit{ef}}
\usepackage{cleveref}
\crefname{originalreference}{original reference}{original references}
\Crefname{originalreference}{Original reference}{Original references}
\makeatletter
\newlabel{lemma-make-pre-equivalence}{{048T}{source}{}{sourcepointer-1}{}}
\newlabel{lemma-make-pre-equivalence@cref}{{[originalreference][0][]048T}{[1][0][]source}{}{}{}}
\makeatother
\crefname{originalreference}{original reference}{original references}
\Crefname{originalreference}{Original reference}{Original references}
\makeatletter
\newlabel{lemma-weak-orbit-pre-equivalence}{{048R}{source}{}{sourcepointer-2}{}}
\newlabel{lemma-weak-orbit-pre-equivalence@cref}{{[originalreference][0][]048R}{[1][0][]source}{}{}{}}
\makeatother
\crefname{originallemma}{original lemma}{original lemmas}
\Crefname{originallemma}{Original lemma}{Original lemmas}
\makeatletter
\newlabel{spaces-morphisms-lemma-base-change-finite-type}{{03XH}{source}{}{sourcepointer-3}{}}
\newlabel{spaces-morphisms-lemma-base-change-finite-type@cref}{{[originallemma][0][]03XH}{[1][0][]source}{}{}{}}
\makeatother
\crefname{originallemma}{original lemma}{original lemmas}
\Crefname{originallemma}{Original lemma}{Original lemmas}
\makeatletter
\newlabel{spaces-morphisms-lemma-large-enough}{{0488}{source}{}{sourcepointer-4}{}}
\newlabel{spaces-morphisms-lemma-large-enough@cref}{{[originallemma][0][]0488}{[1][0][]source}{}{}{}}
\makeatother
\crefname{originallemma}{original lemma}{original lemmas}
\Crefname{originallemma}{Original lemma}{Original lemmas}
\makeatletter
\newlabel{spaces-morphisms-lemma-locally-finite-type-surjective-geometric-points}{{0487}{source}{}{sourcepointer-5}{}}
\newlabel{spaces-morphisms-lemma-locally-finite-type-surjective-geometric-points@cref}{{[originallemma][0][]0487}{[1][0][]source}{}{}{}}
\makeatother
\crefname{originallemma}{original lemma}{original lemmas}
\Crefname{originallemma}{Original lemma}{Original lemmas}
\makeatletter
\newlabel{spaces-morphisms-lemma-permanence-finite-type}{{0462}{source}{}{sourcepointer-6}{}}
\newlabel{spaces-morphisms-lemma-permanence-finite-type@cref}{{[originallemma][0][]0462}{[1][0][]source}{}{}{}}
\makeatother
\title{Weak orbits agree with orbits under locally finite type or large cardinality hypotheses}
\author{The Stacks Project Authors}
\date{stacks-project:groupoids-quotients.tex:local-snapshot-20261012; GNU FDL 1.2 or later}
\begin{document}
\maketitle
\noindent\textit{Editorial note.} All delivered mathematical text is the author's own, copied from the pinned Stacks Project chapter file groupoids-quotients.tex; no statement or proof step was rewritten, reordered or removed.\par\medskip
\noindent\textit{Editorial note.} The counted proof is the author's complete argument: The cardinality bound lambda(R) is the one introduced just before the cited large-enough lemma.\par\medskip
\noindent\textit{Editorial note.} The label lemma-geometric-orbits names the counted statement, and the proof environment immediately following it in the source is delivered verbatim as its proof.\par\medskip
\noindent\textit{Editorial note.} The counted statement is carried by the author's own lemma environment in the chapter source.\par\medskip
\noindent\textit{Allowed standard prerequisites (editorial).}\begin{itemize}
\item Chapter Lemma lemma-make-pre-equivalence (navigation stub, not delivered): replacing a pre-relation R by the pre-equivalence relation R-infinity it generates leaves weak orbits, orbits, and invariant morphisms unchanged, so orbit statements may be proven for pre-equivalence relations
\item Chapter Lemma lemma-weak-orbit-pre-equivalence (navigation stub, not delivered): for a pre-equivalence relation, R-equivalence of k-valued points over extension fields is expressed by the image of R(k) in U(k) times U(k)
\item Morphisms of Spaces, Lemma spaces-morphisms-lemma-permanence-finite-type: local finite type is preserved under composition
\item Morphisms of Spaces, Lemma spaces-morphisms-lemma-base-change-finite-type: local finite type is preserved under base change
\item Morphisms of Spaces, Lemma spaces-morphisms-lemma-locally-finite-type-surjective-geometric-points: a nonempty algebraic space locally of finite type over an algebraically closed field has a k-valued point
\item Morphisms of Spaces, Lemma spaces-morphisms-lemma-large-enough: the same conclusion holds over any algebraically closed field whose cardinality exceeds the bound lambda associated to the morphism
\end{itemize}\medskip
\ifdefined\hypertarget\hypertarget{sourcepointer-1}{}\fi
\noindent\textit{Source reference.} Original reference 048T: sources/stacks-groupoids-quotients/pinned-original-source.tex, source lines 866--866. Prerequisite pointer invoked at 908; it names a standard fact of the same chapter that the counted argument uses, and every step of the counted proof built on it is retained verbatim in the delivered body.\par\medskip
\ifdefined\hypertarget\hypertarget{sourcepointer-2}{}\fi
\noindent\textit{Source reference.} Original reference 048R: sources/stacks-groupoids-quotients/pinned-original-source.tex, source lines 781--781. Prerequisite pointer invoked at 915; it names a standard fact of the same chapter that the counted argument uses, and every step of the counted proof built on it is retained verbatim in the delivered body.\par\medskip
\ifdefined\hypertarget\hypertarget{sourcepointer-3}{}\fi
\noindent\textit{Source reference.} Original lemma 03XH: sources/stacks-chow/tags.txt, source lines 5148--5148. Prerequisite pointer invoked at 926; it names a standard fact of a sibling chapter that the counted argument uses, and every step of the counted proof built on it is retained verbatim in the delivered body.\par\medskip
\ifdefined\hypertarget\hypertarget{sourcepointer-4}{}\fi
\noindent\textit{Source reference.} Original lemma 0488: sources/stacks-chow/tags.txt, source lines 5535--5535. Prerequisite pointer invoked at 937; it names a standard fact of a sibling chapter that the counted argument uses, and every step of the counted proof built on it is retained verbatim in the delivered body.\par\medskip
\ifdefined\hypertarget\hypertarget{sourcepointer-5}{}\fi
\noindent\textit{Source reference.} Original lemma 0487: sources/stacks-chow/tags.txt, source lines 5534--5534. Prerequisite pointer invoked at 929; it names a standard fact of a sibling chapter that the counted argument uses, and every step of the counted proof built on it is retained verbatim in the delivered body.\par\medskip
\ifdefined\hypertarget\hypertarget{sourcepointer-6}{}\fi
\noindent\textit{Source reference.} Original lemma 0462: sources/stacks-chow/tags.txt, source lines 5457--5457. Prerequisite pointer invoked at 922; it names a standard fact of a sibling chapter that the counted argument uses, and every step of the counted proof built on it is retained verbatim in the delivered body.\par\medskip
\begin{lemma}
\label{lemma-geometric-orbits}
Let $S$ be a scheme, and let $B$ be an algebraic space over $S$.
Let $j : R \to U \times_B U$ be a pre-relation over $B$.
Let $\Spec(k) \to B$ be a geometric point of $B$.
\begin{enumerate}
\item If $s, t : R \to U$ are locally of finite type
then weak $R$-equivalence on $U(k)$ agrees with $R$-equivalence, and
weak $R$-orbits agree with $R$-orbits on $U(k)$.
\item If $k$ has sufficiently large cardinality then weak $R$-equivalence
on $U(k)$ agrees with $R$-equivalence, and weak $R$-orbits agree
with $R$-orbits on $U(k)$.
\end{enumerate}
\end{lemma}

\begin{proof}
We first prove (1). Assume $s, t$ locally of finite type. By
Lemma \ref{lemma-make-pre-equivalence}
we may assume that $R$ is a pre-equivalence relation.
Let $k$ be an algebraically closed field over $B$.
Suppose $\overline{u}, \overline{u}' \in U(k)$ are $R$-equivalent.
Then for some extension field $\Omega/k$ there exists
a point $\overline{r} \in R(\Omega)$ mapping to
$(\overline{u}, \overline{u}') \in (U \times_B U)(\Omega)$, see
Lemma \ref{lemma-weak-orbit-pre-equivalence}.
Hence
$$
Z = R \times_{j, U \times_B U, (\overline{u}, \overline{u}')} \Spec(k)
$$
is nonempty. As $s$ is locally of finite type we see that
also $j$ is locally of finite type, see
Morphisms of Spaces, Lemma \ref{spaces-morphisms-lemma-permanence-finite-type}.
This implies $Z$ is a nonempty algebraic space locally of finite type
over the algebraically closed field $k$ (use
Morphisms of Spaces,
Lemma \ref{spaces-morphisms-lemma-base-change-finite-type}).
Thus $Z$ has a $k$-valued point, see
Morphisms of Spaces, Lemma
\ref{spaces-morphisms-lemma-locally-finite-type-surjective-geometric-points}.
Hence we conclude there exists a $\overline{r} \in R(k)$ with
$j(\overline{r}) = (\overline{u}, \overline{u}')$, and we conclude that
$\overline{u}, \overline{u}'$ are $R$-equivalent as desired.

\medskip\noindent
The proof of part (2) is the same, except that it uses
Morphisms of Spaces, Lemma
\ref{spaces-morphisms-lemma-large-enough}
instead of
Morphisms of Spaces, Lemma
\ref{spaces-morphisms-lemma-locally-finite-type-surjective-geometric-points}.
This shows that the assertion holds as soon as $|k| > \lambda(R)$ with
$\lambda(R)$ as introduced just above
Morphisms of Spaces, Lemma
\ref{spaces-morphisms-lemma-locally-finite-type-surjective-geometric-points}.
\end{proof}

\end{document}
